% Additional arithmetic refinement. Counts require the finite-algebra compiler.
\subsection{Removing the scalar preparation coordinate}
\label{subsec:scalar-preparation-zero}

The scalar-height certificate can also use preparation zero, after a
different normalization of its contact parameter and a larger fixed set
of squareclass corrections.  This is a change of arithmetic construction;
it is not an identity on the old full witness space.  The following
argument uses the same explicitly stated arithmetic inputs as the
shared-point definition.

Put
\[
 \delta(T)=T^2+T+23,\qquad \epsilon=\frac{204}{23^{12}},
 \qquad R_\beta=\epsilon^2\frac{r_\beta}{\delta(\beta)^4},
\]
where $\beta=X^2+2$ and $r_\beta$ is the previous scalar contact
polynomial.  Its nonvanishing and $\delta(\beta)>0$ hold for every
rational input.  Define the two fixed prime sets
\[
 \mathcal B=\{5,19,23,349,661,5209\},\qquad
 \mathcal L=\{2,3,11,17,37,67,71\},
\]
and the finite correction set
\[
 \mathcal C'=\left\{\pm\prod_{p\in I}p:
                   I\subseteq\mathcal B\cup\mathcal L\right\}.
\]
It has $2^{14}$ elements.  All these primes belong to the original
exceptional set $S_*$; the change introduces no new exceptional prime
in the height theorem.

Use the same six rational scalar tests and their same orders
$(1,1,1,1,48,48)$, with $J_0=H$.  In their fixed binary norm tree
replace $r_\beta$ by $R_\beta$, including in every leaf denominator.
Thus the six leaves are divided by the assigned powers of $R_\beta$;
the constant leaf remains $1$ on the all-left path.  Denote the root
by $W$, and require $W\ne0$.  Choose $c\in\mathcal C'$ on a finite
branch, and put
\[
 b=cR_\beta W^2,\qquad B=17b,\qquad h=e(B-b^2).
\]
Keep just the two structural coordinates $e,s$, together with the
two finite square roots
\[
 p^2=B(1-e)(1+3e),\qquad
 q^2=\frac{(s^2-Bh)(s^2+3Bh)}B.
\]
The old preparation coordinate $t_*$ and its rational expression
$U_*$ are absent.  The guard $b\ne0$ is retained.

\begin{lemma}[The fixed quadratic and the local correction basis]
\label{lem:scalar-fixed-quadratic}
The polynomial $\delta$ has no root modulo any prime of $\mathcal L$.
The squareclasses of the seven primes in $\mathcal L$ form a basis
of
\[
 \mathbb Q_2^*/\mathbb Q_2^{*2}\ \times\
 \mathbb Q_3^*/\mathbb Q_3^{*2}\ \times\
 \mathbb Q_{17}^*/\mathbb Q_{17}^{*2}.
\]
If an integer $\beta$ makes a prime divide both $r_\beta$ and
$\delta(\beta)$, that prime belongs to $\mathcal B$.  At any finite
scalar contact of a rational $\beta$ outside $S_*$, $\delta(\beta)$
is a unit.
\end{lemma}
\begin{proof}
Modulo two, $\delta=T^2+T+1$.  At the odd primes of $\mathcal L$
its discriminant $-91$ is a quadratic nonresidue.  These are finite
residue computations.  For clarity the local squareclass vectors
are displayed below.  The columns are the valuation at two, the
two unit coordinates for the basis $(-1,5)$ in $\mathbb Z_2^*$,
the valuation and nonsquare-unit coordinate at three, and the
valuation and nonsquare-unit coordinate at seventeen:
\[
\begin{array}{c|rrrrrrr}
2 &1&0&0&0&1&0&0\\
3 &0&1&1&1&0&0&1\\
11&0&1&1&0&1&0&1\\
17&0&0&0&0&1&1&0\\
37&0&0&1&0&0&0&1\\
67&0&1&1&0&0&0&0\\
71&0&1&0&0&1&0&1
\end{array}
\]
Their matrix has rank seven over $\mathbb F_2$, proving the basis
claim by the usual local squareclass descriptions.

The relevant evaluations of the quadratic are
\[
 \begin{split}
 \delta(0)&=23,\qquad \delta(1)=5^2,\\
 \delta(2401)&=5^2\cdot349\cdot661,\\
 32^2\delta(-343/32)&=5^2\cdot5209,\\
 4\delta(1/2)&=5\cdot19.
 \end{split}
\]
For an odd prime dividing the scalar contact polynomial, reduction
therefore gives exactly the asserted list of possible common primes.
The prime two cannot divide $\delta(\beta)$ for integer $\beta$.
The same evaluation argument shows that at a rational finite
contact outside $S_*$ the quadratic is a unit: the contact value
and its fixed denominators are integral there, and its displayed
evaluation is a unit.  This proves both remaining statements.
\end{proof}

\begin{lemma}[Preservation of scalar soundness]
\label{lem:normalized-scalar-soundness}
The modified scalar certificate implies the same estimate
\[
                      h(\beta)\le H\sqrt{|a|}
\]
as the previous scalar-height block.
\end{lemma}
\begin{proof}
First, $\beta=X^2+2$ has no odd contact with infinity at any finite
prime.  Indeed, if $v_p(\beta)<0$, then $v_p(X)<0$ and
$v_p(\beta)=2v_p(X)$ is even.  It therefore suffices to consider
the finite contacts counted by the source's radical $M_5(\beta)$.

At each such odd contact outside $S_*$, the preceding lemma makes
$\delta(\beta)$ a unit.  The fixed factor $\epsilon$ and every
$c\in\mathcal C'$ are also units there.  Consequently
$v_p(R_\beta)=v_p(r_\beta)$ is positive and odd.  The source's
no-cancellation argument for $N_R(u,v)=u^2+Rv^2$ now applies
unchanged, with $R=R_\beta$.  The pole condition on $cRW^2$ forces
the six original rational tests to their assigned orders.  In
particular it gives the order-$48$ divisibility of the scalar index
$a$ at every prime in $M_5(\beta)$.  Thus
$M_5(\beta)^{48}\mid a$, and the original five-contact height
theorem gives the stated estimate, with its original constant $H$.
No test has been divided by an extra, possibly nonunit, normalizer:
only the contact parameter of the norm construction has changed.
\end{proof}

\begin{lemma}[Completeness with preparation zero]
\label{lem:normalized-scalar-completeness}
For every integral input the shared-point construction can satisfy
the modified scalar block together with all of its comparison,
height, and size conditions.
\end{lemma}
\begin{proof}
Fix an integer input, so $\beta\ge2$ is an integer.  First choose
the subset of $\mathcal B$ in $c$ so that
\[
               v_p(cr_\beta)\equiv0\pmod2\quad(p\in\mathcal B).
\]
The factors $\epsilon^2$ and $\delta(\beta)^4$ do not change these
parities.  Choose the sign of $c$ so that $cr_\beta>0$.  Finally
the independent basis $\mathcal L$ supplies a subset making
$cr_\beta$, and hence $cR_\beta$, square in
$\mathbb Q_2,\mathbb Q_3,\mathbb Q_{17}$.  The two prime sets
are disjoint, so this last choice preserves all prescribed parities
at $\mathcal B$.

The reduced denominator of $cR_\beta$ is a square.  At the primes
of $\mathcal B$ its valuation is even.  Away from $\mathcal B$,
if $p\mid\delta(\beta)$, then $p$ divides neither $r_\beta$ nor
$c$: this follows from the resultant calculation and the absence
of roots modulo primes of $\mathcal L$.  Its valuation is then a
multiple of four.  At the remaining primes the valuation is
nonnegative, since the only fixed denominator of $\epsilon$ is
supported at $23\in\mathcal B$.

At two, three, and seventeen the quadratic is a unit and
the factor $204^2=(2^2\cdot3\cdot17)^2$ gives
\[
 v_2(cR_\beta)\ge4,\qquad
 v_3(cR_\beta)\ge2,\qquad
 v_{17}(cR_\beta)\ge2.
\]
The largest absolute value of a correction is
\[
 C_{\max}=\prod_{p\in\mathcal B\cup\mathcal L}p
          =518513917206770054130.
\]
For real $\beta\ge2$ one has
\[
 \frac{|r_\beta|}{\delta(\beta)^4}
 <\frac{1202\cdot204\cdot2}{\beta^3}
 \le61302<2^{16}.
\]
The exact integer inequality
\[
              2^{16}C_{\max}204^2<23^{24}
\]
therefore gives $0<cR_\beta<1$.  All numerical and finite-field
computations in this paragraph are reproduced by the companion
exact checker.

It remains to choose the common point and the anchor point so
that their norm root $W$ is a unit at the finite support of
$cR_\beta$, and $|W-1|<1/2$ in the real absolute value.  The local
condition is the same formal-group construction as before: take
$H_*=[h_0]P$ sufficiently deep at those finitely many places,
then take $S$ sufficiently deeper relative to $H_*$.  Every
nonconstant normalized leaf can consequently have an arbitrarily
large positive valuation.  Even when $v_p(R_\beta)<0$, this makes
each nonconstant summand of every norm node divisible by $p$;
the constant leaf on the all-left path then gives $W\equiv1\pmod p$.

We give the simultaneous real choice explicitly, since it is needed
to remove the preparation coordinate.  Let $G$ be the closure of
the cyclic subgroup generated by the pair of real embeddings of
$P$, let $G^0$ be its identity component, and let $m$ be the order
of $G/G^0$.  The closure of any subgroup $k\mathbb ZP$ contains
$G^0$; if $m\mid k$, it is exactly $G^0$.  The projection of
$G^0$ to either real elliptic curve is nontrivial, since that
projection of $P$ has infinite order.  Thus there are small
nonempty open sets in $G^0$ near its identity which avoid the
identity in both projections.

Choose such a relatively compact open set $\Omega$, so close to
the identity that the two rational coefficients of $f(H)$ are
as small as needed for all $H\in\Omega$.  Shrink it to have
$z(H)$ bounded away from zero at both real embeddings.  Choose
$h_0$ positive, divisible by $64m$, satisfying the local formal
conditions and with $[h_0]P\in\Omega$.  Such a choice follows
from the density of the required finite-index subgroup in $G^0$.
Choose the positive index of $S$ in the required finite-index
subgroup, satisfying its anchor and local conditions, and with
$S$ sufficiently close to the real identity.  Then all six scalar
tests, including the two coefficients of $z(S)/z(H)$, are as small
as prescribed, uniformly for $H\in\Omega$.  At zero values of
all six tests the fixed norm root is exactly one.  Its continuity
therefore gives $|W-1|<1/2$, and in particular $0<W<2$, throughout
this real choice.  The real
conditions on $S$ are compatible with any required size lower
bound because a positive orbit in the finite-index subgroup
returns infinitely often to the chosen open set.

After fixing $S$ and the Lucas root indices, prescribe the common
index modulo $L$, where $64m\mid L$, so that all scalar local
conditions persist near $h_0$ and all small-prime comparison
conditions hold.  Inside $\Omega$ remove the pole loci of the
finitely many comparison functions and then choose a nonempty
open set with compact closure avoiding those loci.  Each locus
has empty interior: its projection is a finite preimage of a
point under a nonzero multiplication map on a real elliptic
curve; all the root labels are nonzero by the completeness choice.
The comparison functions are uniformly bounded on the
chosen compact set.

The simultaneous prime-pattern theorem is now applied as in the
previous sharing proof, including every odd factor of $L$ except
seven in its avoidance modulus.  Its two order moduli $U,V$
satisfy
\[
 \gcd(U,V)=4,\quad \gcd(U,L)=64,\quad \gcd(V,L)=4.
\]
Thus $h\equiv0\pmod U$, $h\equiv4\pmod V$, and
$h\equiv h_0\pmod L$ are compatible.  The final progression
consists of multiples of $m$ and has closure $G^0$ in the real
group.  It therefore contains arbitrarily large positive indices
whose real points lie in the chosen compact-open neighborhood.
The complete comparison and detector construction goes through,
and the scalar inequalities and local units are preserved.

For this final common point,
\[
 b=cR_\beta W^2\quad\hbox{satisfies}\quad
 0<b<4<17.
\]
It has square reduced denominator: at every prime in the support
of $cR_\beta$ the norm root is a unit, and at all other primes
its additional valuation is even.  It is square at two, three,
and seventeen and has valuations there at least $4,2,2$,
respectively.  The preparation-zero pole-completeness lemma
therefore supplies $e,s,p,q$.  All added guards hold, and no
preparation coordinate has been used.
\end{proof}

\begin{theorem}[Conditional forty-three-variable construction]
\label{thm:shared43}
Under the same arithmetic inputs as the shared-point construction,
the scalar modification above preserves the integer input projection.
After finite-algebra compression it uses at most $43$ total rational
coordinates with the indexed finite exclusion, or $45$ with the
semilocal anchor.  The corresponding witness counts are $42$ and
$44$.  These are exact construction bounds, not claims that a
coefficient table of the final polynomial has been expanded.
\end{theorem}
\begin{proof}
Soundness changes only the scalar-height input, which was proved
unchanged.  The preceding lemma proves simultaneous completeness.
The scalar preparation coordinate is deleted, while its two square
roots are retained as finite algebra generators.  Thus the retained
base dimension decreases from $43$ to $42$, or from $45$ to $44$
with the semilocal anchor.  The triangular algebra still has rank
at most $2^{39}$, or $2^{40}$, respectively.  The enlarged fixed
correction set is combined by a product over its branches and
adds no rational coordinate.  The finite-algebra theorem supplies
one final coordinate, giving the asserted totals.
\end{proof}
